\documentclass[a4paper,12pt]{article}
\usepackage[T2A]{fontenc}
\usepackage[utf8]{inputenc}
\usepackage[russian]{babel}
\usepackage{amsmath,amssymb,mathtools}
\renewcommand{\familydefault}{\sfdefault}
\usepackage{geometry}
\geometry{left=20mm,right=20mm,top=20mm,bottom=24mm}
\usepackage{microtype}
\usepackage{tabularx,booktabs,array}
\usepackage{enumitem}
\usepackage{xcolor}
\usepackage{tikz}
\usepackage{fancyhdr}
\usepackage{hyperref}
\definecolor{accent}{HTML}{007B80}
\definecolor{darkaccent}{HTML}{15575A}
\definecolor{ink}{HTML}{2F3335}
\definecolor{alert}{HTML}{B03735}
\definecolor{muted}{HTML}{777777}
\hypersetup{colorlinks=true,linkcolor=darkaccent,pdfauthor={Лёвин Артём Александрович},pdftitle={Ревью домашней работы — Николь Саркисьянц}}
\setlength{\parindent}{0pt}
\setlength{\parskip}{5pt}
\setlist[itemize]{leftmargin=1.6em,itemsep=1pt,topsep=3pt}
\renewcommand{\arraystretch}{1.3}
\pagestyle{fancy}
\fancyhf{}
\fancyhead[L]{\footnotesize\textcolor{muted}{СТЕРЕОМЕТРИЯ · ДОМАШНЯЯ РАБОТА}}
\fancyhead[R]{\footnotesize\textcolor{muted}{17.11.2026}}
\fancyfoot[C]{\fontsize{8}{10}\selectfont\textcolor{muted}{Лёвин Артём Александрович эксклюзивно для Николь Саркисьянц}}
\fancyfoot[R]{\footnotesize\thepage}
\renewcommand{\headrulewidth}{0pt}
\renewcommand{\footrulewidth}{0pt}
\newcommand{\keyidea}{\textcolor{darkaccent}{\textbf{Ключевая идея.}}\,}
\newcommand{\errorlead}{\textcolor{alert}{\textbf{Ошибка.}}\,}
\newcommand{\outcome}{\textcolor{darkaccent}{\textbf{Итог.}}\,}
\newcommand{\parthead}[1]{\par\vspace{4pt}{\bfseries\textcolor{ink}{#1}}\par\vspace{1pt}}
\newcommand{\taskhead}[2]{\phantomsection\label{task:#1}\section*{Задание №\,#1}\textcolor{muted}{#2}\par\smallskip}
\begin{document}
\begin{titlepage}
\thispagestyle{empty}
\centering
\vspace*{30mm}
{\large\textcolor{darkaccent}{МАТЕМАТИКА}\par}
\vspace{29mm}
{\LARGE\textcolor{ink}{Ревью домашней работы}\par}
\vspace{9mm}
{\Large\textcolor{accent}{Призмы и пирамиды}\par}
\vspace{16mm}
{\normalsize\textcolor{muted}{Проверка решений по стереометрии}\par}
\vfill
{\large Николь Саркисьянц\par}
\vspace{11mm}
{17 ноября 2026 г.\par}
\vspace{17mm}
{\small Лёвин Артём Александрович\par}
\vspace{2mm}
{\small эксклюзивно для Николь Саркисьянц\par}
\vspace{14mm}
\begin{tikzpicture}
\draw[accent,line width=.65pt] (0,0) .. controls (1.4,.65) and (2.4,-.55) .. (4,0)
 .. controls (5.4,.55) and (6.5,-.48) .. (8,0)
 .. controls (9.1,.35) and (10,-.25) .. (11,0);
\end{tikzpicture}
\end{titlepage}

\section*{Сводная таблица}
\textcolor{muted}{В таблицу включены только задания, требующие исправления или завершения. Нажмите на номер задания, чтобы перейти к подробному разбору.}
\par\vspace{5mm}
\small
\begin{tabularx}{\linewidth}{@{}>{\centering\arraybackslash}p{10mm} >{\raggedright\arraybackslash}p{25mm} >{\raggedright\arraybackslash}X >{\raggedright\arraybackslash}p{23mm} >{\raggedright\arraybackslash}p{31mm}@{}}
\toprule
\textbf{№} & \textbf{Тема} & \textbf{Суть ошибки} & \textbf{Ваш ответ} & \textbf{Правильный ответ}\\
\midrule
\hyperref[task:5]{\textbf{5}} & Объём призмы & Коэффициент изменения высоты взят неверно & $81$ & $54$\\
\addlinespace[4pt]
\hyperref[task:7]{\textbf{7}} & Полная поверхность пирамиды & У пирамиды ошибочно учтены два основания; решение не завершено & Не указан & $36+9\sqrt3$\\
\addlinespace[4pt]
\hyperref[task:10]{\textbf{10}} & Апофема пирамиды & Высота пирамиды подставлена вместо апофемы; периметр основания неверен & $S_{\text{бок}}=36$; апофема не найдена & $l=5$; $S_{\text{бок}}=60$\\
\bottomrule
\end{tabularx}
\normalsize

\newpage
\thispagestyle{fancy}
\vspace*{\fill}
\begin{center}
{\LARGE\textcolor{darkaccent}{Разбор ошибок}}\\[12pt]
{\normalsize\textcolor{muted}{Пошаговое исправление и самостоятельное закрепление}}
\end{center}
\vspace*{\fill}
\newpage

\taskhead{5}{Тема: как меняется объём правильной треугольной призмы}
\parthead{Условие}
Объём правильной треугольной призмы равен $18$. Сторону основания увеличили \textbf{втрое}, а высоту уменьшили \textbf{втрое}. Найдите новый объём.

\parthead{Ваше решение и ответ}
Использована формула объёма через площадь равностороннего треугольника. Далее в записи появляется отношение
\[
\frac{18}{V_{\text{нов}}}=\frac{2}{9},
\qquad V_{\text{нов}}=\frac{18\cdot9}{2}=81.
\]
В промежуточной записи изменение высоты также рассматривается как уменьшение в два раза.

\parthead{Где возникает ошибка}
\errorlead Формула объёма и идея сравнивать коэффициенты изменения выбраны правильно. Площадь равностороннего треугольника при увеличении его стороны втрое возрастает в $9$ раз. \textbf{Первое неверное действие} --- использование множителя $2$ для изменения высоты. По условию высоту уменьшили \textbf{втрое}, следовательно, новая высота равна одной трети первоначальной.

\parthead{Почему это неверно}
Объём прямой призмы вычисляется умножением площади основания на высоту. Если площадь возросла в $9$ раз, а высота уменьшилась в $3$ раза, общий коэффициент изменения равен $9:3=3$. Число $2$ в отношении $2/9$ не соответствует условию, поэтому получен завышенный объём $81$.

\parthead{Ключевая идея}
\keyidea Для равностороннего треугольника $S=\frac{a^2\sqrt3}{4}$. При изменении стороны в $k$ раз площадь меняется в $k^2$ раз, а объём призмы дополнительно зависит от изменения высоты.

\parthead{Исправление от правильного шага}
\[
S_{\text{нов}}=9S_{\text{стар}},\qquad
H_{\text{нов}}=\frac{H_{\text{стар}}}{3}.
\]
Тогда
\[
\frac{V_{\text{нов}}}{V_{\text{стар}}}
=9\cdot\frac13=3.
\]

\parthead{Полное правильное решение}
Пусть $a$ и $H$ --- первоначальные сторона основания и высота. Тогда
\[
V_{\text{стар}}=\frac{a^2\sqrt3}{4}\,H=18.
\]
После изменения размеров получаем
\begin{align*}
V_{\text{нов}}
&=\frac{(3a)^2\sqrt3}{4}\cdot\frac{H}{3}\\
&=\frac{9a^2\sqrt3}{4}\cdot\frac{H}{3}\\
&=3\cdot\frac{a^2\sqrt3}{4}\,H\\
&=3\cdot18=54.
\end{align*}

\parthead{Проверка}
Новая площадь основания больше в девять раз, высота меньше в три раза. Следовательно, объём должен быть втрое больше первоначального: $18\cdot3=54$. Это совпадает с расчётом.

\outcome \textbf{$V_{\text{нов}}=54$.}

\parthead{Задача на закрепление}
Объём правильной треугольной призмы равен $24$. Сторону основания увеличили вдвое, а высоту уменьшили в четыре раза. Найдите новый объём.

\newpage
\taskhead{7}{Тема: площадь полной поверхности правильной треугольной пирамиды}
\parthead{Условие}
В правильной треугольной пирамиде сторона основания равна $6$, боковое ребро равно $5$. Найдите площадь полной поверхности.

\parthead{Ваше решение и ответ}
В записи использована формула
\[
S_{\text{полн}}=S_{\text{бок}}+2S_{\text{осн}}.
\]
При этом площадь основания найдена правильно:
\[
S_{\text{осн}}=\frac{6^2\sqrt3}{4}=9\sqrt3.
\]
Численное значение полной поверхности не получено.

\parthead{Где возникает ошибка}
\errorlead Первый неверный шаг --- прибавление \textbf{двух} площадей основания. У пирамиды только \textbf{одно основание}; остальные грани --- боковые треугольники. Формула с двумя основаниями относится к призме. Правильная формула здесь:
\[
S_{\text{полн}}=S_{\text{бок}}+S_{\text{осн}}.
\]
\textbf{Сохранить можно} правильно найденную площадь равностороннего треугольника $9\sqrt3$.

\parthead{Ключевая идея}
\keyidea У правильной треугольной пирамиды три равных боковых грани. Каждая из них --- равнобедренный треугольник со сторонами $5$, $5$ и $6$. Высоту боковой грани находят по теореме Пифагора.

\parthead{Исправление и полное правильное решение}
\textbf{Шаг 1.} Найдём площадь основания:
\[
S_{\text{осн}}=\frac{6^2\sqrt3}{4}=9\sqrt3.
\]
\textbf{Шаг 2.} Рассмотрим одну боковую грань. Её высота $l$ делит основание длины $6$ на два отрезка по $3$. В полученном прямоугольном треугольнике гипотенуза равна боковому ребру $5$:
\[
l^2+3^2=5^2,\qquad l^2=25-9=16,
\qquad l=4.
\]
\textbf{Шаг 3.} Площадь одной боковой грани:
\[
S_1=\frac12\cdot6\cdot4=12.
\]
\textbf{Шаг 4.} Боковых граней три, поэтому
\[
S_{\text{бок}}=3\cdot12=36.
\]
\textbf{Шаг 5.} Прибавим площадь \emph{одного} основания:
\[
S_{\text{полн}}=36+9\sqrt3.
\]

\parthead{Проверка}
Три боковые грани дают площадь $36$. Основание имеет положительную площадь $9\sqrt3$, поэтому полная поверхность больше боковой ровно на площадь одного основания. Приближённо $S_{\text{полн}}\approx51{,}6$.

\outcome \textbf{$S_{\text{полн}}=36+9\sqrt3$.}

\parthead{Задача на закрепление}
В правильной треугольной пирамиде сторона основания равна $8$, а боковое ребро --- $5$. Найдите площадь полной поверхности.

\newpage
\taskhead{10}{Тема: апофема и площадь боковой поверхности правильной четырёхугольной пирамиды}
\parthead{Условие}
В правильной четырёхугольной пирамиде сторона основания равна $6$, высота пирамиды равна $4$. Найдите высоту боковой грани и площадь боковой поверхности.

\parthead{Ваше решение и ответ}
Записаны данные $a=6$, $H=4$. Высота боковой грани не найдена. Для боковой поверхности записано:
\[
S_{\text{бок}}=\frac12\cdot4\cdot18=36.
\]

\parthead{Где возникает ошибка}
\errorlead Правильно выделены исходные величины, но \textbf{первый неверный шаг} --- непосредственная подстановка высоты пирамиды $4$ в вычисление боковой площади. Высота пирамиды $H$ перпендикулярна плоскости основания; высота боковой грани $l$ проходит по наклонной треугольной грани. В этой задаче это разные отрезки.

Второй недочёт: число $18$ нельзя принять за периметр квадратного основания со стороной $6$. Периметр квадрата равен $4\cdot6=24$. Кроме того, в ответе отсутствует первая запрошенная величина --- высота боковой грани.

\parthead{Ключевая идея}
\keyidea Апофема правильной пирамиды, её вертикальная высота и половина стороны квадратного основания образуют прямоугольный треугольник. По теореме Пифагора:
\[
l^2=H^2+\left(\frac a2\right)^2.
\]

\parthead{Исправление от правильного шага}
Из записанных данных $a=6$ и $H=4$ сначала вычислим апофему:
\[
l=\sqrt{4^2+3^2}=5.
\]
Только после этого используем формулу $S_{\text{бок}}=\frac12Pl$ с периметром квадрата $P=24$.

\parthead{Полное правильное решение}
\textbf{Шаг 1.} Обозначим вершину пирамиды через $S$, центр квадратного основания --- через $O$, середину одной стороны --- через $M$. В правильной пирамиде $SO$ --- высота, $SM$ --- апофема. Так как $SO$ перпендикулярна основанию, треугольник $SOM$ прямоугольный.

\textbf{Шаг 2.} Расстояние от центра квадрата до середины стороны:
\[
OM=\frac a2=\frac62=3.
\]
\textbf{Шаг 3.} По теореме Пифагора:
\[
SM^2=SO^2+OM^2=4^2+3^2=25,
\qquad SM=5.
\]
\textbf{Шаг 4.} Основание боковой треугольной грани равно $6$, высота этой грани равна $5$. Площадь одной грани:
\[
S_1=\frac12\cdot6\cdot5=15.
\]
\textbf{Шаг 5.} Таких граней четыре:
\[
S_{\text{бок}}=4\cdot15=60.
\]

\parthead{Проверка}
Тем же результатом обладает расчёт через периметр:
\[
S_{\text{бок}}=\frac12Pl=\frac12\cdot24\cdot5=60.
\]
Оба способа совпали. Апофема $5$ больше вертикальной высоты $4$, что согласуется с прямоугольным треугольником.

\outcome \textbf{Высота боковой грани $5$; площадь боковой поверхности $60$.}

\parthead{Задача на закрепление}
В правильной четырёхугольной пирамиде сторона основания равна $10$, высота пирамиды равна $12$. Найдите высоту боковой грани и площадь боковой поверхности.

\end{document}